\scanpage{179}
\mspage{36 bis}

\inserted{2.20.} Contr'exemple à la proposition fausse 2.19.

Soit $X$ un\deleted{\unsure{e catég}} objet de Cat, rajoutons lui \deleted{un objet} la catégorie finale $e$ comme fermé, i.e. considérons un objet $X'$ de Cat dans lequel $X$ soit un sous objet ouvert, $e$ le fermé complémentaire :

\diagram{\begin{tikzcd}[ampersand replacement=\&]
X\arrow[r,hook,"i"] \& X' \& e\arrow[l,hook',"j"']
\end{tikzcd}}

On sait (cf $\underline{\overline{\mathrm{XII}}}$) qu'une telle opération est déterminée, à isom. près, par la donnée d'un objet

\[
H\in\operatorname{Ob}X^\wedge,
\]

en posant

\[
H(x)=\operatorname{Hom}_{X'}(x,e),
\]

la catégorie $X_{/e}$ s'identifiant (à isom près) à $X_{/H}$. Elle est donc $0$-connexe ssi l'objet $H$ de $X^\wedge$ est $0$-connexe. Cela signifie aussi que le morphisme $i$ est cocofinal, i.e. que $i^{\unsure{o|0}}:X^o\hookrightarrow X'^{o}$ est cofinal. Supposons maintenant $X$ $\mathbb W_\infty$-asphérique et tot. $0$-connexe \deleted{($W$ étant un loc. fond. satisfaisant $W(\unsure{\mathrm{As}|\lim})$)}, i.e l'actuel Loc(6), cf $\underline{\overline{\unsure{\mathrm{XVI}|\mathrm{XVII}}}}$.) La proposition 2.19 affirme que sous ces conditions, $X'$ (qui \unsure{est} lui aussi tot $0$-connexe) est $=$ $W$-asphérique. Il s'ensuivrait \deleted{que} donc que $i\in\mathbb W_\infty$. Mais par $\underline{\overline{\mathrm{XII}}}$ th. 4 cor. \unsure{4|5} (p.~149) il s'ensuivrait que

% archive sheet number: 179
